\documentclass[12pt]{article}

\usepackage[T2A]{fontenc}
\usepackage[utf8]{inputenc}
\usepackage[russian,english]{babel}
\usepackage{array}
\usepackage{enumitem}
\usepackage{geometry}
\usepackage{graphicx}
\usepackage{hyperref}
\usepackage{makecell}
\usepackage{tabularx}
\usepackage{tikz}
\usepackage{xltabular}
\usepackage[table]{xcolor}
\usepackage{listings}

\geometry{
    a4paper,
    left=1.5cm,
    right=1.5cm,
    top=1.5cm,
    bottom=2.5cm
}

\hypersetup{hidelinks}
\setlength{\parindent}{0cm}
\setlength{\parskip}{0.35\baselineskip}
\setlength{\emergencystretch}{1.5em}
\setlength{\tabcolsep}{7pt}
\renewcommand{\arraystretch}{1.12}
\setcellgapes{2pt}
\makegapedcells
\keepXColumns

\definecolor{commitred}{RGB}{236, 48, 48}
\definecolor{commitblue}{RGB}{52, 83, 255}
\newcolumntype{L}[1]{>{\raggedright\arraybackslash}p{#1}}
\newcolumntype{Y}{>{\raggedright\arraybackslash}X}
\newcommand{\cmdcell}[1]{\parbox[t]{\linewidth}{\ttfamily\small\raggedright\strut #1\strut}}

\definecolor{codebg}{RGB}{248,248,248}
\definecolor{commentcolor}{RGB}{0,128,0}
\definecolor{keywordcolor}{RGB}{0,0,128}
\definecolor{stringcolor}{RGB}{163,21,21}
\lstset{
    backgroundcolor=\color{codebg},
    basicstyle=\ttfamily\small,
    breaklines=true,
    frame=single,
    rulecolor=\color{black!30},
    keywordstyle=\color{keywordcolor},
    commentstyle=\color{commentcolor},
    stringstyle=\color{stringcolor},
    numbers=left,
    numberstyle=\tiny\color{gray},
    stepnumber=1,
    numbersep=8pt,
    captionpos=b,
    showstringspaces=false,
    tabsize=2
}

\begin{document}

\begin{titlepage}
    \centering
    Федеральное государственное автономное \\
    образовательное учреждение высшего образования \\
    \textbf{<<Национальный исследовательский университет ИТМО>>} \\

    \vfill
    Веб-программирование \\
    \vspace{0.5\baselineskip}
    \textbf{Лабораторная работа №3} \\
    \vspace{0.5\baselineskip}
    Вариант 1403 \\

    \vfill

    \flushright
    Выполнил: \\
    Студент группы P3214 \\

    \vspace{4\baselineskip}

    \centering
    г. Санкт-Петербург, 2026

\end{titlepage}

\clearpage
\tableofcontents
\newpage

\section{Цель работы}

Изучить и реализовать сценарий сборки проекта с помощью Apache Ant, обеспечивающий компиляцию, тестирование, генерацию документации и упаковку Java-проекта из лабораторной работы №3 по дисциплине «Веб-программирование».

\section{Техническое задание}

Разработать Ant-сценарий, удовлетворяющий следующим требованиям:

\begin{enumerate}[leftmargin=*, label=\arabic*)]
    \item Все переменные и константы должны быть вынесены в отдельный файл параметров.
    \item \texttt{MANIFEST.MF} должен содержать информацию о версии и запускаемом классе.
    \item Цель \texttt{compile} должна осуществлять компиляцию исходных кодов.
    \item Цель \texttt{build} должна вызывать \texttt{compile} и упаковывать приложение в JAR-архив.
    \item Цель \texttt{clean} должна удалять скомпилированные классы и временные файлы.
    \item Цель \texttt{test} должна запускать JUnit-тесты проекта после сборки.
    \item Цель \texttt{doc} должна добавлять MD5 и SHA-1 в \texttt{MANIFEST.MF}, генерировать Javadoc и добавлять его в архив.
    \item Цель \texttt{team} должна получать предыдущие git-ревизии, собирать их и упаковывать в ZIP через \texttt{build}.
\end{enumerate}

\section{Исходные данные}

Проект расположен в каталоге \texttt{opi3}. Исходная структура проекта включает:

\begin{itemize}[leftmargin=1.5em]
    \item \texttt{src/main/java} -- исходные Java-классы;
    \item \texttt{src/test/java} -- тестовые классы;
    \item \texttt{src/main/resources} -- ресурсы и конфигурационные файлы;
    \item \texttt{build.properties} -- файл параметров;
    \item \texttt{build.xml} -- Ant-сценарий.
\end{itemize}

\section{Ход работы}

\subsection{Подготовка параметров}

Создан файл \texttt{build.properties}, содержащий все настройки проекта. Это позволяет изменять пути, версии и параметры компиляции без изменения основного сценария.

\subsection{Реализация цели compile}

Цель \texttt{compile} реализована отдельным блоком Ant-сценария. Она выполняет:

\begin{itemize}[leftmargin=1.5em]
    \item создание каталогов сборки;
    \item компиляцию файлов из \texttt{src/main/java} в \texttt{build/classes};
    \item копирование ресурсов из \texttt{src/main/resources} в каталог классов.
\end{itemize}

\subsection{Реализация цели build}

Цель \texttt{build} зависит от \texttt{compile} и выполняет:

\begin{itemize}[leftmargin=1.5em]
    \item генерацию файла \texttt{MANIFEST.MF};
    \item упаковку классов и ресурсов в JAR-архив \texttt{dist/web3.jar};
    \item включение информации о версии, основном классе и дате сборки.
\end{itemize}

\subsection{Реализация цели clean}

Цель \texttt{clean} удаляет каталоги \texttt{build/} и \texttt{dist/}, а также временные файлы типа \texttt{*.bak} и \texttt{*~}. Это возвращает проект в исходное состояние.

\subsection{Реализация цели test}

Цель \texttt{test} запускается после \texttt{build} и выполняет следующие действия:

\begin{itemize}[leftmargin=1.5em]
    \item компиляцию тестовых классов в \texttt{build/test-classes};
    \item копирование тестовых ресурсов;
    \item запуск JUnit-тестов;
    \item сохранение результатов в \texttt{build/reports}.
\end{itemize}

\subsection{Реализация цели doc}

Цель \texttt{doc} реализует:

\begin{itemize}[leftmargin=1.5em]
    \item генерацию документации Javadoc по всем классам;
    \item расчёт MD5 и SHA-1 для итогового JAR;
    \item обновление \texttt{MANIFEST.MF} с контрольными суммами;
    \item пересоздание архива с учетом Javadoc.
\end{itemize}

\subsection{Реализация цели team}

Цель \texttt{team} выполняет:

\begin{itemize}[leftmargin=1.5em]
    \item проверку наличия git-репозитория;
    \item получение двух предыдущих ревизий;
    \item сборку каждой ревизии через цель \texttt{build};
    \item упаковку полученных JAR-файлов в \texttt{dist/team-builds.zip}.
\end{itemize}

\section{Результаты}

В результате работы были созданы:

\begin{itemize}[leftmargin=1.5em]
    \item \texttt{build.properties} -- конфигурация параметров Ant;
    \item \texttt{build.xml} -- сценарий Ant с целями \texttt{compile}, \texttt{build}, \texttt{clean}, \texttt{test}, \texttt{doc}, \texttt{team};
    \item \texttt{report/title.tex} -- титульная страница отчёта;
    \item \texttt{report/report.tex} -- отчёт по заданной структуре.
\end{itemize}

\subsection{Список целей}

\begin{xltabular}{\textwidth}{|L{0.25\textwidth}|Y|}
\hline
\rowcolor{white}
\textbf{Цель} & \textbf{Назначение} \\
\hline
compile & компиляция исходных кодов проекта \\
\hline
build & создание JAR и вызов \texttt{compile} \\
\hline
clean & удаление артефактов сборки \\
\hline
test & запуск JUnit-тестов после сборки \\
\hline
doc & генерация документации и запись контрольных сумм \\
\hline
team & сборка двух предыдущих ревизий и архивирование \\
\hline
\end{xltabular}

\section{Приложение}

\subsection{Листинг файла build.properties}

\lstinputlisting[caption={Конфигурация параметров Ant}]{../build.properties}

\subsection{Листинг файла build.xml}

\lstinputlisting[caption={Ант-сценарий сборки}]{../build.xml}

\section{Вывод}

Разработанный Ant-сценарий полностью соответствует техническому заданию. Он обеспечивает разделение конфигурации и логики, воспроизводимую сборку, автоматический запуск тестов, генерацию документации и контрольных сумм, а также сборку предыдущих версий через git.

Файл \texttt{MANIFEST.MF} содержит информацию о версии и запускаемом классе, а также добавленные контрольные суммы MD5 и SHA-1 после цели \texttt{doc}. Отчёт подготовлен в стиле ЛР2 и содержит подробное описание выполнения работы.

\end{document}
